% BHU PYQ -- Manually transcribed & error-corrected
% Paper: BCH-321 : Financial Analysis
% Exam: B.Com. (Hons.) Semester VI Examination, 2022-23

\documentclass[11pt,a4paper]{article}
\usepackage[a4paper,top=2.5cm,bottom=2.5cm,left=2.8cm,right=2.8cm,headheight=14pt]{geometry}
\usepackage{amsmath,amssymb}
\usepackage[T1]{fontenc}
\usepackage[utf8]{inputenc}
\usepackage{lmodern,microtype,fancyhdr,enumitem,hyperref,lastpage,parskip}

\hypersetup{
    colorlinks=true,
    urlcolor=black,
    linkcolor=black,
    pdftitle={BCH-321: Financial Analysis -- BHU B.Com. (Hons.) Sem VI 2022-23}
}

\pagestyle{fancy}
\fancyhf{}
\renewcommand{\headrulewidth}{0.4pt}
\renewcommand{\footrulewidth}{0.4pt}
\fancyhead[L]{\small   \textit{Banaras Hindu University}}
\fancyhead[R]{\small   \textit{BCH-321: Financial Analysis}}
\fancyfoot[C]{\small Page  \thepage\ of \pageref{LastPage}}


\newlist{parts}{enumerate}{2}
\setlist[parts,1]{label=(\alph*),leftmargin=*,itemsep=5pt,topsep=3pt}
\setlist[parts,2]{label=(\roman*),leftmargin=*,itemsep=3pt,topsep=2pt}

\newcommand{\pts}[1]{\hfill   \textit{[#1]}}


\begin{document}


\begin{center}
  {\large\bfseries BANARAS HINDU UNIVERSITY}\\[2pt]
  {\normalsize B.Com. (Hons.) --- Semester VI Examination, 2022-23}\\[6pt]
  \rule{0.8\linewidth}{0.5pt}\\[6pt]
  {\large\bfseries Paper No. BCH-321: Financial Analysis}\\[4pt]
  {\normalsize Time: 3 Hours\hfill Maximum Marks: 70}
\end{center}

\rule{\linewidth}{0.4pt}

\smallskip



\noindent   \textit{Instructions: Attempt all questions. All questions carry equal marks (14 marks each).}

\bigskip




\noindent   \textbf{Question 1.}
``Inability to understand and deal with financial data proves to be a serious handicap in managing a business successfully.'' Explain the statement.\pts{14}

\centerline{   \textbf{OR}}

``For the purposes of financial statement analysis, the items of income statement and balance sheet cannot be taken as they are. They are required to be reclassified so as to unfold and understand the hidden information.'' Explain the statement.\pts{14}

\medskip\rule{\linewidth}{0.2pt}




\noindent   \textbf{Question 2.}
From the following information, prepare the balance sheet of A Ltd:\pts{14}

\begin{center}

\begin{tabular}{ll}
  \hline
  Particulars & Ratio / Amount \\
  \hline
  Net Working Capital & Rs. 12,00,000 \\
  Cost of Sales to Closing Stock Ratio & 8 \\
  Gross Profit Ratio & 20\% \\
  Fixed Assets Turnover Ratio (on cost of sales) & 2 \\
  Current Ratio & 3 \\
  Liquid Ratio & 2 \\
  Debtors Collection Period & 2 months \\
  Ratio of Fixed Assets to Shareholders' Net Worth & 0.75 \\
  Ratio of Reserves and Surplus to Capital & 0.60 \\
  \hline
\end{tabular}
\end{center}

\centerline{   \textbf{OR}}

What is ratio analysis? Why is it done? What are its limitations?\pts{14}

\medskip\rule{\linewidth}{0.2pt}


\newpage




\noindent   \textbf{Question 3.}
From the following details, you are required to prepare the Cash Flow Statement of Alfa Ltd:\pts{14}


\begin{center}
   \textbf{Statement of Profit and Loss for the year ended 31st March, 2022}
\vskip 5pt

\begin{tabular}{lll}
  \hline
  Particulars & Note No. & Amount (Rs.) \\
  \hline
   \textbf{I. Revenue} & & \\
  \quad Revenue from Operations & 1 & 32,00,000 \\
  \quad Other Income & 2 & 50,000 \\
   \textbf{Total Revenue} & &   \textbf{32,50,000} \\
  \hline
   \textbf{II. Expenses} & & \\
  \quad Cost of Goods Sold & & 20,00,000 \\
  \quad Finance Cost & 3 & 15,000 \\
  \quad Depreciation & & 2,10,000 \\
  \quad Other Expenses & 4 & 7,91,000 \\
   \textbf{Total Expenses} & &   \textbf{30,16,000} \\
  \hline
   \textbf{Profit Before Tax (I - II)} & &   \textbf{2,34,000} \\
  \quad Less: Provision for Income Tax & & 92,000 \\
  \hline
   \textbf{Profit for the year} & &   \textbf{1,42,000} \\
  \hline
\end{tabular}
\end{center}


\begin{center}
   \textbf{Notes to Accounts:}
\vskip 5pt

\begin{tabular}{lll}
  \hline
  Particulars & & Amount (Rs.) \\
  \hline
   \textbf{1. Revenue from Operations} & & \\
  \quad Sales of goods in trade & & 32,00,000 \\
  \hline
   \textbf{2. Other Income} & & \\
  \quad Government compensation for loss in riots & & 50,000 \\
  \hline
   \textbf{3. Finance Cost} & & \\
  \quad Interest on Debentures & & 15,000 \\
  \hline
   \textbf{4. Other Expenses} & & \\
  \quad Operating Expenses & & 7,90,000 \\
  \quad Cost of issue of debentures written off & & 1,000 \\
   \textbf{Total} & &   \textbf{7,91,000} \\
  \hline
\end{tabular}
\end{center}


\newpage

\begin{center}
   \textbf{Comparative Balance Sheet Items:}
\vskip 5pt

\begin{tabular}{lll}
  \hline
  Particulars & Year 2021 (Rs.) & Year 2022 (Rs.) \\
  \hline
  Inventories & 1,80,000 & 2,20,000 \\
  Trade Debtors & 40,000 & 38,000 \\
  Bills Receivables & 30,000 & 55,000 \\
  Balance with Bank and Cash in Hand & 1,17,000 & 2,48,000 \\
  Trade Creditors & 78,000 & 95,000 \\
  Bills Payables & 20,000 & 15,000 \\
  Outstanding Expenses & 31,000 & 44,000 \\
  \hline
\end{tabular}
\end{center}

   \textbf{Additional Information:}
The following important transactions have taken place during the year ended 31st March, 2022:

\begin{parts}
  \item Fully paid equity shares of the face value of Rs. 2,00,000 were allotted at a premium of 20\%.
  \item 10\% debentures for Rs. 3,00,000 were redeemed at a premium of 2\%.
  \item Land was purchased for Rs. 1,50,000 and the consideration was discharged by the allotment to the vendor of zero per cent convertible debentures for the amount.
  \item Dividend and dividend distribution tax thereon for the year ended 31st March, 2021 amounting to Rs. 1,15,000 was paid.
  \item Income tax paid during the year totaled Rs. 95,000.
\end{parts}

\centerline{   \textbf{OR}}

What is Funds? What is flow of funds? How is fund from operation calculated?\pts{14}

\medskip\rule{\linewidth}{0.2pt}




\noindent   \textbf{Question 4.}
What is financial forecasting? Why is it done?\pts{14}

\centerline{   \textbf{OR}}

How is a projected income statement prepared? Explain by giving a suitable example.\pts{14}

\medskip\rule{\linewidth}{0.2pt}




\noindent   \textbf{Question 5.}
The following particulars are supplied to you to calculate the value of goodwill on the basis of capitalisation of super profit:\pts{14}

\begin{parts}
  \item Capital Invested: Rs. 10,00,000
  \item Trading Results:
  
\begin{center}
  
\begin{tabular}{ll}
    \hline
    Year & Profit (Rs.) \\
    \hline
    2015 & 50,000 \\
    2016 & 1,50,000 \\
    2017 & 3,00,000 \\
    2018 & 5,00,000 \\
    2019 & 8,00,000 \\
    \hline
  \end{tabular}
  \end{center}
  \item Standard rate of return in the industry is 20\%.
\end{parts}
Also throw light on as to how the trend in trading results could be taken into account while calculating the super profit.

\centerline{   \textbf{OR}}

Explain a method of share valuation which you think is the best.\pts{14}

\vfill
\rule{\linewidth}{0.4pt}

\begin{center}\small
\href{https://bhu-exam.vercel.app/}{Click here to visit our website}\\[4pt]
\href{https://me-aryan.vercel.app/}{Click here to visit our contributor}\\[4pt]
\href{https://quashan-me.vercel.app/}{Click here to visit our contributor}
\end{center}

\end{document}
